\section{v1.3.0: declared numeric input and reference bands}
Existing numeric entries keep their point-decimal syntax. The new helper uses
explicit rules and assigns a canonical literal locally; every call resets its
input defaults. The output is data, independent of the display formatter.
\begin{demo}{A declared decimal comma}
\gradbarsparsenumber[decimal=comma,grouping=point]
 {1.234,5}{\purchaseAmount}
\gradbar[max=2000,width=55mm,precision=1]{\purchaseAmount}
\end{demo}
\key{decimal=point|comma} and \key{grouping=none|point|comma|space} must differ.
Grouped integers need a first group of 1--3 digits, then groups of exactly three;
fractions and exponents are ungrouped. Space means ASCII space. Leading/trailing
spaces are trimmed. Empty input and uppercase NA mean missing; zero is numeric.
Scientific notation uses the declared decimal marker and an ungrouped exponent.
The existing digit and magnitude bounds still apply; expressions, units, NaN
and infinity remain errors. No ambiguous separator is guessed.
\subsection{CSV input rules are column declarations}
The file below contains quoted decimal-comma fields. Only explicit CSV scalar
columns (number, bar, bullet, heat) accept input rules. Text and the loaded dataset
remain literal; canonical values feed selection, scales, ranks and summaries.
\begin{demo}{Library purchases; amount is a raw numeric total}
\gradbarsloadcsv{purchases}{docs/data/regional-input.csv}
\begin{center}
\begin{gradtable}[data=purchases,style=regional]
\gradcolumn[source=Item,width=26mm]{item}{Item}
\gradcolumn[source=Amount,type=number,width=30mm,
 input decimal=comma,input grouping=point,
 aggregate=sum,options={precision=2}]{amount}{Amount}
\gradsummary{Total}
\end{gradtable}
\end{center}
\end{demo}
Table \key{style=regional} sets accessible colors, spacing and decimal number
alignment, without selecting an input locale. Display formatting is independent.
\clearpage
\subsection{A trend and a constant reference band}
\key{\gradsparkband} takes a series, lower bound and upper bound. A fixed range
or named scale and line type are mandatory. Bounds must increase within that
range. \key{band color} and \key{band opacity} set the fill; opacity is in [0,1].
\begin{demo}{Circulation rates; fixed reference interval 45--65 percent}
\begingroup
\gradscale{circulation}{min=0,max=100}
\gradbarssetup{width=55mm,height=5ex,
 precision=0,unit={\%},label width=12mm}
\begin{tabular}{lc}
Collection & Circulation; reference: 45--65\%\\
\hline
Literature & \gradsparkband[scale group=circulation]
 {35,55,NA,60,68}{45}{65}\\
Science & \gradsparkband[scale group=circulation]
 {60,60,60,60,60}{45}{65}\\
History & \gradsparkband[scale group=circulation]
 {NA,NA,NA,NA,NA}{45}{65}\\
\end{tabular}
\endgroup
\end{demo}
The band is a supplied reference, not observations or a computed confidence
interval. It never changes the range or fills missing positions. NA breaks the
trend; all-missing series retain the band. Print mode uses gray with dashed
boundaries. Automatic alternatives retain the series, range and reference in
one Figure. Existing spark commands and their unsupported-option errors remain.
Scalar custom backends need explicit TikZ fallback for \key{spark-band}.

Standalone sources: \texttt{docs/examples/regional-input.tex} and its Chinese
counterpart. Full input matrix: \texttt{docs/regional-input.md}.
